\documentclass[11pt]{amsart}
\usepackage[a4paper,margin=1in]{geometry}
\usepackage{amsmath,amssymb,amsthm,mathtools}
\usepackage[T1]{fontenc}
\usepackage{lmodern}
\usepackage{microtype}
\usepackage{xcolor}
\usepackage{hyperref}
\usepackage{enumitem}
\hypersetup{colorlinks=true,linkcolor=blue!50!black,citecolor=blue!50!black,urlcolor=blue!50!black}

\newtheorem{theorem}{Theorem}
\newtheorem{proposition}[theorem]{Proposition}
\newtheorem{lemma}[theorem]{Lemma}
\newtheorem{claim}[theorem]{Claim}

\newcommand{\E}{\mathbb{E}}
\newcommand{\optprior}{\mathrm{OPT}_{\mathrm{prior}}}
\newcommand{\optpost}{\mathrm{OPT}_{\mathrm{post}}}

\title{A $5/4$ Lower-Bound Family for the A Priori / A Posteriori Gap}
\author{}
\date{}

\begin{document}
\maketitle

We consider the following family of instances. Let $n\ge 5$ be odd, let $v_1,\dots,v_n$ be the vertices of $K_n$ (indices modulo $n$), and fix the Hamiltonian cycle
\[
C:=v_1v_2\cdots v_nv_1.
\]
For each $i\in[n]$, subdivide the edge $v_iv_{i+1}$ by a new vertex $x_i$, and let $d$ be the shortest-path metric of the resulting graph. We take $r:=v_1$ as depot. The client set is
\[
V:=\{v_2,\dots,v_n,x_1,\dots,x_n\}.
\]
Every original client $v_2,\dots,v_n$ is active with probability $1$, and each subdividing vertex $x_i$ is active independently with probability $1/2$.

The main goal of this note is to show that this instance yields the asymptotically tight lower bound $5/4$ for the ratio between the optimal a priori and a posteriori values.

\begin{theorem}\label{thm:main}
For every $\varepsilon>0$, there exists a stochastic TSP instance such that
\[
\frac{\optprior}{\optpost}\ge \frac54-\varepsilon.
\]
More precisely, for every odd $n\ge 5$, the single-cycle subdivision instance above satisfies
\[
\optprior=\frac{15}{8}n
\qquad\text{and}\qquad
\optpost\le \frac32 n+3.
\]
Consequently,
\[
\frac{\optprior}{\optpost}
\ge
\frac{(15/8)n}{(3/2)n+3}
\xrightarrow[n\to\infty]{} \frac54.
\]
\end{theorem}

The proof is split into two parts: first the exact a priori value, then the a posteriori upper bound.

\section*{1. The a priori value}

We first record the metric structure. For original vertices,
\[
d(v_i,v_{i+1})=2,
\qquad
 d(v_i,v_j)=1\quad\text{if } j\not\equiv i\pm 1 \pmod n.
\]
For subdividing vertices,
\[
d(x_i,v_i)=d(x_i,v_{i+1})=1,
\qquad
 d(x_i,v_j)=2\quad\text{if } j\notin\{i,i+1\}.
\]
For two subdividing vertices,
\[
d(x_i,x_j)=2
\quad\Longleftrightarrow\quad
x_i\text{ and }x_j\text{ correspond to adjacent edges of }C,
\]
and otherwise $d(x_i,x_j)=3$.

\begin{proposition}\label{prop:prior}
For every odd $n\ge 5$, the above instance satisfies
\[
\optprior=\frac{15}{8}n.
\]
\end{proposition}

\begin{proof}
Let $T$ be an arbitrary master tour. Since the original vertices are always active, deleting all subdividing vertices from $T$ yields a cyclic ordering
\[
(r=u_1,u_2,\dots,u_n,r)
\]
of the original vertices. For each $j\in[n]$, let
\[
Y_j=(y_{j,1},\dots,y_{j,k_j})
\]
be the ordered sequence of subdividing vertices that appears between $u_j$ and $u_{j+1}$ in $T$, where $u_{n+1}:=r$. Then
\[
\sum_{j=1}^n k_j=n.
\]

For an ordered pair of original vertices $u,v$ and an ordered sequence $Y=(y_1,\dots,y_k)$ of distinct subdividing vertices, let
\[
\Phi(u;Y;v)
\]
denote the expected length of the realized path from $u$ to $v$ obtained from the segment
\[
(u,y_1,\dots,y_k,v)
\]
after deleting inactive subdividing vertices. By linearity of expectation,
\[
\E[c(T|_A)]
=
\sum_{j=1}^n \Phi(u_j;Y_j;u_{j+1}).
\]
So it suffices to lower-bound each segment.

\medskip
\noindent\textbf{Step 1: a convenient exact formula.}
Set
\[
z_0:=u,\qquad z_i:=y_i\ (1\le i\le k),\qquad z_{k+1}:=v.
\]
For $0\le i<j\le k+1$, let $E_{ij}$ be the event that $z_i$ and $z_j$ are retained and all vertices strictly between them are deleted. Then $z_i$ and $z_j$ become consecutive after shortcutting if and only if $E_{ij}$ occurs, hence
\[
\Phi(u;Y;v)=\sum_{0\le i<j\le k+1} d(z_i,z_j)\Pr(E_{ij}).
\]
Since each subdividing vertex is retained independently with probability $1/2$,
\[
\Pr(E_{0,k+1})=2^{-k},\qquad
\Pr(E_{0,t})=2^{-t},\qquad
\Pr(E_{t,k+1})=2^{-(k-t+1)},\qquad
\Pr(E_{ij})=2^{-(j-i+1)}.
\]
Therefore
\begin{align*}
\Phi(u;Y;v)
&=
2^{-k}d(u,v)
+\sum_{t=1}^k 2^{-t}d(u,y_t)
+\sum_{t=1}^k 2^{-(k-t+1)}d(y_t,v) \\
&\qquad
+\sum_{1\le i<j\le k}2^{-(j-i+1)}d(y_i,y_j).
\end{align*}

Introduce the indicators
\[
\delta:=\mathbf 1_{\{u,v\text{ consecutive on }C\}},
\qquad
\alpha_t:=\mathbf 1_{\{y_t\text{ is incident with }u\}},
\qquad
\beta_t:=\mathbf 1_{\{y_t\text{ is incident with }v\}},
\]
and, for $1\le i<j\le k$,
\[
\gamma_{ij}:=\mathbf 1_{\{y_i,y_j\text{ correspond to adjacent edges of }C\}}.
\]
Then
\[
d(u,v)=1+\delta,
\qquad
 d(u,y_t)=2-\alpha_t,
\qquad
 d(y_t,v)=2-\beta_t,
\qquad
 d(y_i,y_j)=3-\gamma_{ij}.
\]
Substituting and simplifying the geometric sums gives
\[
\Phi(u;Y;v)=1+\frac{3k}{2}+2^{-k}\delta-D(Y),
\]
where
\[
D(Y):=
\sum_{t=1}^k 2^{-t}\alpha_t
+\sum_{t=1}^k 2^{-(k-t+1)}\beta_t
+\sum_{1\le i<j\le k}2^{-(j-i+1)}\gamma_{ij}.
\]

\medskip
\noindent\textbf{Step 2: lower bounds for short segments.}
We claim that every segment with $k:=|Y|$ satisfies
\[
\Phi(u;Y;v)\ge
\begin{cases}
1, & k=0,\\[1mm]
2, & k=1,\\[1mm]
\dfrac{11}{4}, & k=2,\\[2mm]
\dfrac{29}{8}, & k=3.
\end{cases}
\]
Moreover, equality for $k=2$ is attained by $(v_i,x_i,x_{i+1},v_{i+1})$, and equality for $k=3$ is attained by $(v_{i+1},x_{i+1},x_i,x_{i-1},v_i)$.

For $k=0$, we simply have $\Phi(u;\emptyset;v)=d(u,v)\ge 1$.

For $k=1$, say $Y=(y)$, the exact formula gives
\[
\Phi(u;y;v)=\frac12 d(u,v)+\frac12\bigl(d(u,y)+d(y,v)\bigr).
\]
If $u$ and $v$ are consecutive on $C$, then $d(u,v)=2$ and $d(u,y)+d(y,v)\ge 2$. If they are not consecutive, then $d(u,v)=1$ and no subdividing vertex is incident with both endpoints, so $d(u,y)+d(y,v)\ge 3$. In either case, $\Phi(u;y;v)\ge 2$.

For $k=2$, the formula becomes
\[
\Phi(u;y_1,y_2;v)=4+\frac\delta4-D(Y).
\]
If $\delta=0$, then each $y_t$ is incident with at most one of $u,v$, so $\alpha_t+\beta_t\le 1$ for $t=1,2$. Hence
\[
D(Y)
=\frac12\alpha_1+\frac14\alpha_2+\frac14\beta_1+\frac12\beta_2+\frac14\gamma_{12}
\le 1+\frac14=\frac54,
\]
which yields $\Phi(u;y_1,y_2;v)\ge 11/4$. If $\delta=1$, then at most one of $y_1,y_2$ can be the subdivider of the edge $uv$; therefore the endpoint contribution is at most $5/4$, and again $\gamma_{12}\le 1$, so $D(Y)\le 3/2$. Hence
\[
\Phi(u;y_1,y_2;v)\ge 4+\frac14-\frac32=\frac{11}{4}.
\]
Equality is attained by $(v_i,x_i,x_{i+1},v_{i+1})$.

For $k=3$, the formula becomes
\[
\Phi(u;y_1,y_2,y_3;v)=\frac{11}{2}+\frac\delta8-D(Y),
\]
where
\begin{align*}
D(Y)
&=\frac12\alpha_1+\frac14\alpha_2+\frac18\alpha_3
+\frac18\beta_1+\frac14\beta_2+\frac12\beta_3 \\
&\qquad +\frac14\gamma_{12}+\frac18\gamma_{13}+\frac14\gamma_{23}.
\end{align*}
If $\delta=0$, then each $y_t$ is incident with at most one of $u,v$, so the endpoint part is at most
\[
\frac12+\frac14+\frac12=\frac54,
\]
and among three cycle edges there are at most two adjacent pairs, whence
\[
\frac14\gamma_{12}+\frac18\gamma_{13}+\frac14\gamma_{23}\le \frac12.
\]
Thus $D(Y)\le 7/4$, and therefore
\[
\Phi(u;y_1,y_2,y_3;v)\ge \frac{11}{2}-\frac74=\frac{15}{4}>\frac{29}{8}.
\]
If $\delta=1$, then at most one of $y_1,y_2,y_3$ can be incident with both endpoints. Hence the endpoint part is at most $3/2$, while the adjacency part is still at most $1/2$. Thus $D(Y)\le 2$, and so
\[
\Phi(u;y_1,y_2,y_3;v)\ge \frac{11}{2}+\frac18-2=\frac{29}{8}.
\]
Equality is attained by $(v_{i+1},x_{i+1},x_i,x_{i-1},v_i)$.

\medskip
\noindent\textbf{Step 3: long segments.}
Now suppose $k\ge 4$. Since an original vertex is incident with exactly two subdividing vertices,
\[
\sum_{t=1}^k 2^{-t}\alpha_t\le 2^{-1}+2^{-2}=\frac34,
\qquad
\sum_{t=1}^k 2^{-(k-t+1)}\beta_t\le \frac34.
\]
Moreover, the graph on $\{y_1,\dots,y_k\}$ in which two vertices are adjacent when the corresponding cycle edges are adjacent has maximum degree at most $2$, hence at most $k$ edges. Since every coefficient $2^{-(j-i+1)}\le 1/4$,
\[
\sum_{1\le i<j\le k}2^{-(j-i+1)}\gamma_{ij}\le \frac{k}{4}.
\]
Therefore
\[
D(Y)\le \frac34+\frac34+\frac{k}{4}=\frac32+\frac{k}{4},
\]
and so
\[
\Phi(u;Y;v)
\ge
1+\frac{3k}{2}-\left(\frac32+\frac{k}{4}\right)
=
\frac54k-\frac12.
\]
Since $k\ge 4$,
\[
\frac54k-\frac12\ge 1+\frac{7k}{8}.
\]
Hence every segment with $k\ge 1$ satisfies
\[
\Phi(u;Y;v)\ge 1+\frac{7k}{8},
\]
while the case $k=0$ gives $\Phi(u;\emptyset;v)\ge 1$.

\medskip
\noindent\textbf{Step 4: summing over segments.}
Using $\sum_j k_j=n$,
\begin{align*}
\E[c(T|_A)]
&=
\sum_{j=1}^n \Phi(u_j;Y_j;u_{j+1}) \\
&\ge
\sum_{j:k_j=0} 1 + \sum_{j:k_j\ge 1}\left(1+\frac{7k_j}{8}\right) \\
&=
 n + \frac78\sum_{j=1}^n k_j \\
&=
 n+\frac78n
=
\frac{15}{8}n.
\end{align*}
Since $T$ was arbitrary, this proves $\optprior\ge 15n/8$.

\medskip
\noindent\textbf{Step 5: matching upper bound for odd $n$.}
Write $n=2m+1$ with $m\ge 2$. Consider the master tour
\begin{align*}
T^\star:=(&v_n,v_2,x_2,x_1,x_n,v_1,
  v_4,x_4,x_3,v_3,
  v_6,x_6,x_5,v_5,
  \dots, \\
 &v_{2m},x_{2m},x_{2m-1},v_{2m-1},v_n).
\end{align*}
This tour contains exactly one $3$-segment,
\[
(v_2,x_2,x_1,x_n,v_1),
\]
whose expected contribution is $29/8$, and exactly $m-1$ $2$-segments,
\[
(v_{2t},x_{2t},x_{2t-1},v_{2t-1})
\qquad (2\le t\le m),
\]
whose expected contributions are all $11/4$. The remaining $m+1$ segments are direct edges between original vertices, each of length $1$. Therefore
\begin{align*}
\E[c(T^\star|_A)]
&=
\frac{29}{8}+(m-1)\cdot \frac{11}{4}+(m+1) \\
&=
\frac{30m+15}{8}
=
\frac{15(2m+1)}{8}
=
\frac{15}{8}n.
\end{align*}
Hence $\optprior\le 15n/8$, and the proof is complete.
\end{proof}

\section*{2. The a posteriori value}

We now turn to the realized deterministic optimum. For a realization $A$, let
\[
S(A):=\{i\in[n]: x_i\text{ is active in }A\}
\]
and write $m:=|S(A)|$.

\begin{proposition}\label{prop:post}
For every realization $A$ of the single-cycle subdivision instance,
\[
\mathrm{OPT}_A\le n+|S(A)|+3.
\]
Consequently,
\[
\optpost\le \frac32 n+3.
\]
\end{proposition}

\begin{proof}
Fix a realization $A$, and set $S:=S(A)$ and $m:=|S|$.

If $m=n$, then every subdivider is active, and the full subdivided cycle
\[
v_1,x_1,v_2,x_2,\dots,v_n,x_n,v_1
\]
is a feasible tour of cost $2n=n+m\le n+m+3$. So we may assume $m<n$.

Contract every active edge $v_iv_{i+1}$ of the base cycle $C$; equivalently, contract every maximal active run on $C$ into a single block. The resulting quotient is again a cycle, now on
\[
q:=n-m
\]
blocks. Each block is either a single original vertex or a path of the form
\[
v_a,x_a,v_{a+1},x_{a+1},\dots,x_b,v_{b+1}
\]
coming from a maximal consecutive active run. Traversing such a block from one end to the other costs exactly $2\ell$, where $\ell$ is the number of active subdividers in the block, and summing over all blocks gives total internal cost exactly $2m$.

Now consider only the movement \emph{between} blocks. Label the $q$ blocks cyclically by $w_1,\dots,w_q$. If two blocks are consecutive on this quotient cycle, then the distance between the corresponding boundary vertices is at most $2$. If they are not consecutive, then the corresponding boundary originals are nonconsecutive on the original cycle $C$, so their distance is exactly $1$. Thus the inter-block metric is dominated by the standard $q$-vertex cycle metric in which consecutive pairs have distance $2$ and all other pairs have distance $1$.

Let $\tau(q)$ denote the optimal TSP value in that standard $q$-vertex cycle metric. Then the discussion above shows that
\[
\mathrm{OPT}_A\le 2m+\tau(q).
\]
We now bound $\tau(q)$. For $q\ge 5$, there is a Hamiltonian cycle using only nonconsecutive pairs, hence $\tau(q)=q$. Indeed:
\begin{itemize}[leftmargin=2em]
    \item if $q$ is odd, the step-$2$ cycle
    \[
    1,3,5,\dots,q,2,4,\dots,q-1,1
    \]
    has total cost $q$;
    \item if $q$ is even, the cycle
    \[
    1,3,5,\dots,q-1,2,q,q-2,\dots,4,1
    \]
    also has total cost $q$.
\end{itemize}
For the small cases, one checks directly that
\[
\tau(1)\le 2,\qquad \tau(2)=4,\qquad \tau(3)=6,\qquad \tau(4)=6.
\]
In particular,
\[
\tau(q)\le q+3 \qquad \text{for every } q\ge 1.
\]
Therefore
\[
\mathrm{OPT}_A\le 2m+(q+3)=2m+(n-m)+3=n+m+3.
\]
Since $m=|S(A)|$, this proves the first claim.

Taking expectation and using $\E[|S(A)|]=n/2$, we obtain
\[
\optpost
=\E_A[\mathrm{OPT}_A]
\le
\E_A[n+|S(A)|+3]
=
 n+\frac n2+3
=
\frac32 n+3.
\]
\end{proof}

\section*{3. Proof of the ratio}

\begin{proof}[Proof of Theorem~\ref{thm:main}]
For odd $n\ge 5$, Proposition~\ref{prop:prior} gives
\[
\optprior=\frac{15}{8}n,
\]
while Proposition~\ref{prop:post} gives
\[
\optpost\le \frac32 n+3.
\]
Hence
\[
\frac{\optprior}{\optpost}
\ge
\frac{(15/8)n}{(3/2)n+3}.
\]
The right-hand side tends to $5/4$ as $n\to\infty$ through odd integers. Therefore, given any $\varepsilon>0$, we may choose odd $n$ large enough that
\[
\frac{(15/8)n}{(3/2)n+3}\ge \frac54-\varepsilon.
\]
For that choice of $n$, the single-cycle subdivision instance satisfies
\[
\frac{\optprior}{\optpost}\ge \frac54-\varepsilon.
\]
This proves the theorem.
\end{proof}

\end{document}
